\usepackage{iftex}
\ifPDFTeX
	% Use T1 font encoding for better hyphenation, accented characters
	\usepackage[T1]{fontenc}
	% Allows direct use of Unicode characters
	\usepackage[utf8]{inputenc}
\else
	% Use modern unicode fonts
	\usepackage{fontspec}
	%\setmainfont{TeX Gyre Termes}
	%\setmonofont{DejaVu Sans Mono}
\fi

% Use Latin Modern fonts instead of the default fonts
\usepackage{lmodern}
% Improves typography/readability
\usepackage{microtype}

% make the PDF files generated by pdflatex "searchable and copyable"
%\usepackage{cmap}

% Define custom colors
% - table: include rowcolor for table bg coloring
% - dvipsnames: include more default color names
\usepackage[table,dvipsnames]{xcolor}

% Advanced math formatting/symbols
\usepackage{amsmath, amssymb}

% Provides customizable lists: itemize, enumerate
\usepackage{enumitem}

% PDF metadata and document links (must be before glossaries package)
% Customize later e.g. for metadata using
% \hypersetup{
% 	pdfauthor={...},
% 	pdftitle={...},
% 	pdfsubject={...},
% 	pdfkeywords={...}
% 	% Example config to colorize all links/anchors blue
% 	hidelinks=false,
% 	colorlinks=true,
% 	linkcolor=blue,
% 	urlcolor=blue,
% 	citecolor=blue,
% }
\usepackage[%
	% Adds hyperlinks but removes the colored boxes/underlines around links
	hidelinks,
	% Adds pdf metadata
	pdftex,
]{hyperref}
% Automatic detection what is being referenced (table, figure, ...) -> \cref{...}
% \cref{fig:trace_others,fig:ping_others}, \cref{fig:trace_others}
% Customize names for other languages using:
% \crefname{section}{Abschnitt}{Abschnitte}
% \Crefname{section}{Abschnitt}{Abschnitte}
% \crefname{figure}{Abbildung}{Abbildungen}
% \Crefname{figure}{Abbildung}{Abbildungen}
% \crefname{table}{Tabelle}{Tabellen}
% \Crefname{table}{Tabelle}{Tabellen}
% \crefname{equation}{Gleichung}{Gleichungen}
% \Crefname{equation}{Gleichung}{Gleichungen}
\usepackage{cleveref}

% Automatically expand used acronyms
% \gls{}           → prints full form first time, acronym afterwards
% \glspl{}         → prints plural
% \glsentryshort{} → print acronym only (e.g. for captions and section titles)
% \acrshort{}      → acronym only
% \acrfull{}       → full form only
\usepackage[acronym]{glossaries}
\glsdisablehyper % disable links from acronyms to acronym list
\makeglossaries % create index

% Bibliography/Citations
\usepackage[
	backend=biber,
	% Style (style=numeric,style=ieee)
	style=ieee,
	% Sorting (sorting=none, sorting=ynt)
	sorting=none,
]{biblatex}

% Typeset units and numbers
\usepackage{siunitx}

% Force figures to not move to the next page: \begin{figure}[H]
\usepackage{float}

% Include graphics/images: \includegraphics -> image
\usepackage{graphicx}
% SVG support: \includesvg{...} -> converts to image so no text selection or vector graphics
%\usepackage{svg}
% Better looking tables (\toprule, \midrule, \bottomrule)
\usepackage{booktabs}
% Customize captions for figures/tables
\usepackage{caption}
% Subfigures inside a figure
\usepackage{subcaption}

% Compute total page count
\usepackage{lastpage}

% Modern Syntax highlighting for code
\usepackage{minted}

% Embed CSV data as tables
\usepackage{csvsimple}
% Easy multiline cells
% > \makecell{\textbf{Broadcast-} \\ \textbf{Adresse}}
\usepackage{makecell}
% Break weird lines for tables: tabular=|c|c|>{\RaggedRight\arraybackslash}p{7cm}|
\usepackage{ragged2e}

% For tables with footnotes -> \begin{tabular}...\tnote{2}...\end{tabular}\begin{tablenotes} \item[1] X \item[2] Y \end{tablenotes}
\usepackage{threeparttable}

% KOMA package for headers and footers
\usepackage{scrlayer-scrpage}

% Quotations: \enquote{text} -> "text"
\usepackage{csquotes}

% Create advanced macros/commands: \NewDocumentCommand{name}{[options]}{code}
\usepackage{xparse}

% Add TODO notes for drafts: \todo{text}
\usepackage{todonotes}

% Allow landscape pages: \begin{landscape}...\end{landscape}
\usepackage{pdflscape}

% Visualize keyboard shortcuts: \keys{Ctrl + C}
\usepackage{menukeys}
